A retinopatia diabética é uma complicação do diabetes que pode evoluir de forma silenciosa e causar danos irreversíveis à visão quando não identificada e tratada a tempo. Ela é uma das principais causas de cegueira evitável em pessoas com diabetes no Brasil, o que torna seu diagnóstico precoce fundamental. Por meio do exame de fundoscopia, é possível identificar indicadores relevantes, tais como exsudatos moles e duros, que auxiliam na avaliação da doença. Contudo, em diversas regiões o acesso a especialistas e à infraestrutura adequada para o diagnóstico precoce ainda é limitado, além da subjetividade na interpretação do exame, que varia até entre especialistas e é maior quando realizado por profissionais não especializados. Este trabalho consiste na realização de um estudo comparativo entre seis arquiteturas de redes de aprendizado profundo, representativas de três famílias, utilizando modelos tradicionais, como VGG16 e ResNet50, modelos leves, como MobileNetV3 e EfficientNet-B0, e modelos baseados em Transformers, como ViT e Swin Transformer, de forma a identificar os modelos mais adequados para a segmentação de exsudatos duros e moles e para a classificação do grau da retinopatia diabética. Na segmentação, as seis redes serão usadas como codificadores de uma mesma U-Net, com treino no IDRiD e teste também no DDR, e na classificação elas serão treinadas para dar o grau da doença de 0 a 4 com as bases APTOS e EyePACS, sendo testadas também no Messidor-2 sem nenhum ajuste. Todos os modelos serão treinados sob o mesmo protocolo, com validação cruzada, e comparados por testes estatísticos com correção para várias comparações e pelo custo computacional, medido pelo número de parâmetros e de operações, pelo tempo de inferência e pela memória. Os modelos escolhidos serão utilizados na construção de uma solução de diagnóstico auxiliado por computador (CAD) que mostre as lesões, o grau previsto e um alerta de edema macular, de forma a auxiliar, com baixo custo computacional, no diagnóstico de retinopatia diabética em regiões com recursos médicos limitados.

% Separe as palavras-chave por ponto e vírgula
\palavraschave{retinopatia diabética; segmentação e classificação de imagens; aprendizado profundo; exsudatos; diagnóstico auxiliado por computador.}
